\section{Method}
\label{sec:method}

\subsection{Two samples behind one ensemble}
\label{sec:method-setup}

There are \(K\) frozen language models, indexed by \(i=1,\ldots,K\).
Problems are drawn independently from an unknown distribution \(\mathcal D\);
the labeled training benchmark is \(q_1,\ldots,q_Q\).  Each problem \(q\) has a
finite candidate set \(\mathcal A(q)\), with
\(|\mathcal A(q)|\le A\), and gold answer \(a^\star(q)\).
Model \(i\) induces an answer distribution \(p_i(\cdot\mid q)\).  We do not
observe that distribution.  Instead, we draw \(N\) independent generations
\(Y_{i,1}(q),\ldots,Y_{i,N}(q)\) and form
\begin{equation}
\widehat p_i(a\mid q)
=
\frac1N\sum_{n=1}^N\mathbf 1\{Y_{i,n}(q)=a\}.
\label{eq:phat}
\end{equation}
Thus \(\widehat p_i(a\mid q)\) is simply a normalized vote count.  The vectors
\(p_i(\cdot\mid q)\) and \(\widehat p_i(\cdot\mid q)\) lie in the probability
simplex over \(\mathcal A(q)\): they have one nonnegative coordinate per answer
and their coordinates sum to one.  This construction exposes two independent
sample sizes.  \(Q\) controls learning across problems; \(N\) controls how well
each model's answer distribution is measured on a problem.

An ensemble uses a convex weight vector
\(w\in\Delta_K=\{w\in\mathbb R_+^K:\sum_iw_i=1\}\).  It assigns answer \(a\)
the population and empirical scores
\begin{equation}
s_q(a;w)=\sum_{i=1}^K w_i p_i(a\mid q),
\qquad
\widehat s_q(a;w)=\sum_{i=1}^K w_i\widehat p_i(a\mid q).
\label{eq:scores}
\end{equation}
The corresponding predictors return the highest-scoring answer.  Equal
weights recover ordinary voting; nonuniform weights give a weighted plurality
vote.  In the large-\(N\) limit, empirical voting converges to
\(\arg\max_a s_q(a;w)\).

\subsection{The hardest contender determines correctness}
\label{sec:method-margin}

For each wrong answer \(a\ne a^\star(q)\), define the \(K\)-vector
\(\Delta_{q,a}=p_\cdot(a^\star(q)\mid q)-p_\cdot(a\mid q)\), and define
\(\widehat\Delta_{q,a}\) analogously.  The inner product
\(\langle w,\Delta_{q,a}\rangle\) is the ensemble score of the gold answer
minus the score of \(a\).  It is positive when the gold answer beats that
contender and negative when the contender wins.  Because a correct prediction
must beat every wrong answer, the relevant margin is the smallest such gap:
\begin{equation}
M_q(w)=\min_{a\ne a^\star(q)}\langle w,\Delta_{q,a}\rangle,
\qquad
\widehat M_q(w)=
\min_{a\ne a^\star(q)}
\langle w,\widehat\Delta_{q,a}\rangle.
\label{eq:margins}
\end{equation}
Taking the minimum is conservative, not favorable: it selects the hardest
wrong answer.  Hence \(M_q(w)>0\) means that every gap is positive and the
ensemble is correct; \(M_q(w)\le0\) means that some wrong answer ties or wins.
We count ties as errors and write
\begin{equation}
\ell_q(w)=\mathbf 1\{M_q(w)\le0\},
\qquad
R(w)=\mathbb P_{q\sim\mathcal D}\{M_q(w)\le0\}.
\label{eq:risk}
\end{equation}

For a concrete example, consider two models and two wrong answers with
\(\Delta_{q,a_1}=(0.4,-0.1)\), \(\Delta_{q,a_2}=(0.1,0.3)\), and
\(w=(0.5,0.5)\).  The two gaps are \(0.15\) and \(0.20\), so the margin is
\(0.15>0\): the ensemble is correct, and \(a_1\) is the binding contender.

The same definition has a useful geometric reading.  With one wrong answer,
the correct weights form one halfspace of the \((K-1)\)-dimensional weight
simplex.  With several wrong answers, correctness requires all corresponding
halfspace inequalities.  Its closure is therefore a polytope, and the strict
correct region is the relative interior selected by those inequalities.
Across \(Q\) problems, at most \(Q(A-1)\) such hyperplanes partition the
simplex.  This is why the statistical complexity depends primarily on the
weight dimension \(K-1\), with only logarithmic dependence on the number of
contenders in the finite-sample arrangement bound.  Figure~\ref{fig:weight-geometry}
restores the geometric view developed in the reviewed draft.

\begin{figure}[t]
\centering
\begin{tikzpicture}[scale=1.9,every node/.style={font=\scriptsize}]
  \begin{scope}
    \coordinate (A) at (0,0);
    \coordinate (B) at (2.4,0.3);
    \coordinate (C) at (1.1,1.0);
    \coordinate (D) at (0.9,2.2);
    \draw[gray,dashed] (A)--(C) (B)--(C) (C)--(D);
    \draw[thick] (A)--(B) (A)--(D) (B)--(D);
    \coordinate (P1) at (0.405,0.99);
    \coordinate (P2) at (1.725,1.155);
    \coordinate (P3) at (1.01,1.54);
    \fill[blue!14,opacity=0.8] (P1)--(P2)--(D)--cycle;
    \fill[blue!14,opacity=0.8] (P1)--(P3)--(D)--cycle;
    \fill[blue!14,opacity=0.8] (P2)--(P3)--(D)--cycle;
    \draw[blue!70!black,thick] (P1)--(P2);
    \draw[blue!70!black,thick,dashed] (P2)--(P3)--(P1);
    \node[blue!55!black] at (0.95,1.78) {correct};
    \node[below left] at (A) {$e_1$};
    \node[below right] at (B) {$e_2$};
    \node[right] at (C) {$e_3$};
    \node[above] at (D) {$e_4$};
    \node[align=center] at (1.1,-0.55)
      {(a) One wrong answer: one plane\\[-1pt]cuts the simplex, and the correct\\[-1pt]weights form one halfspace.};
  \end{scope}
  \begin{scope}[xshift=4.3cm]
    \coordinate (A) at (0,0);
    \coordinate (B) at (2.4,0.3);
    \coordinate (C) at (1.1,1.0);
    \coordinate (D) at (0.9,2.2);
    \draw[gray,dashed] (A)--(C) (B)--(C) (C)--(D);
    \draw[thick] (A)--(B) (A)--(D) (B)--(D);
    \coordinate (Q1) at (0.80,0.72);
    \coordinate (Q2) at (1.42,0.86);
    \coordinate (Q3) at (1.30,1.40);
    \coordinate (Q4) at (0.70,1.18);
    \fill[blue!16,opacity=0.85] (Q1)--(Q2)--(Q3)--(Q4)--cycle;
    \draw[blue!70!black,thick] ($(Q1)!-0.9!(Q2)$)--($(Q2)!-0.9!(Q1)$);
    \draw[blue!70!black,thick] ($(Q2)!-0.6!(Q3)$)--($(Q3)!-0.9!(Q2)$);
    \draw[blue!70!black,thick] ($(Q3)!-0.9!(Q4)$)--($(Q4)!-0.9!(Q3)$);
    \draw[blue!70!black,thick] ($(Q4)!-0.6!(Q1)$)--($(Q1)!-0.9!(Q4)$);
    \node[blue!55!black] at (1.05,1.05) {$P$};
    \node[below left] at (A) {$e_1$};
    \node[below right] at (B) {$e_2$};
    \node[right] at (C) {$e_3$};
    \node[above] at (D) {$e_4$};
    \node[align=center] at (1.1,-0.55)
      {(b) Several wrong answers: each adds\\[-1pt]one plane, and their halfspaces\\[-1pt]intersect in the region $P$.};
  \end{scope}
\end{tikzpicture}
\caption{Correct-weight geometry for one question with \(K=4\).  The weight
simplex is a tetrahedron whose vertices are pure model choices.  Each contender
creates a decision hyperplane.  With one contender, correctness selects one
side of one plane; with several contenders, all inequalities must hold, giving
the relative interior of a polytope.}
\label{fig:weight-geometry}
\end{figure}

\subsection{Buffered empirical risk minimization}
\label{sec:method-buffer}

Raw empirical risk minimization can reward a weight vector for exploiting
finite-\(N\) sign errors near zero.  We instead require the empirical gold
score to clear a positive buffer.  For \(\varepsilon\ge0\), define
\begin{equation}
\widehat R_\varepsilon(w)
=
\frac1Q\sum_{j=1}^Q
\mathbf 1\{\widehat M_{q_j}(w)\le\varepsilon\},
\qquad
\widehat w_\varepsilon\in
\arg\min_{w\in\Delta_K}\widehat R_\varepsilon(w).
\label{eq:buffered-erm}
\end{equation}
The theory chooses
\(\varepsilon=\varepsilon_N(\delta)
=\sqrt{2\log(8KAQ/\delta)/N}\), the uniform training-sample scale of the
finite-generation error.  The constant is not essential; the statistical
role is.  A problem is counted as empirically solved only when the observed
gold score beats every observed contender by more than the uncertainty
radius.

\begin{algorithm}[t]
\caption{Margin-buffered weighted voting}
\label{alg:buffered}
\begin{algorithmic}[1]
\STATE \textbf{Input:} labeled problems \(q_{1:Q}\), models \(1{:}K\),
generation count \(N\), confidence \(\delta\)
\FOR{each training problem \(q_j\) and model \(i\)}
  \STATE draw \(N\) answers and compute the frequencies in
  Eq.~(\ref{eq:phat})
\ENDFOR
\STATE set
\(\varepsilon_N\leftarrow\sqrt{2\log(8KAQ/\delta)/N}\)
\STATE solve Eq.~(\ref{eq:buffered-erm}) for
\(\widehat w_{\varepsilon_N}\)
\STATE \textbf{return} the learned weights \(\widehat w_{\varepsilon_N}\)
\STATE \textbf{Deploy:} estimate fresh frequencies
\(\widehat p_i(\cdot\mid q)\)
\STATE \textbf{predict}
\(\widehat f_N(q)=\arg\max_a\sum_i\widehat w_{\varepsilon_N,i}
\widehat p_i(a\mid q)\)
\end{algorithmic}
\end{algorithm}

The main theorem isolates the effect of finite generations during weight
fitting: it controls the clean population risk
\(R(\widehat w_{\varepsilon_N})\), whose predictor uses the population scores
\(s_q(a;\widehat w_{\varepsilon_N})\).  The practical
\(\widehat f_N\) above incurs an additional prediction-time perturbation from
fresh generations.  Appendix~\ref{app:legacy-support} gives a sign-reliability
bound for any fixed weight; extending it uniformly to the data-dependent
\(\widehat w_{\varepsilon_N}\) requires boundary control beyond the one-sided
comparator condition used by Theorem~\ref{thm:buffered-main}.

\subsection{Exact optimization as a MILP}
\label{sec:method-milp}

The buffered objective is discontinuous but admits a direct mixed-integer
linear implementation.  Introduce \(z_j\in\{0,1\}\), where \(z_j=1\) claims
that problem \(j\) clears the buffer.  A solver uses a small strictness
tolerance \(\tau>0\) and any
\(B_{\rm M}\ge 1+\varepsilon_N+\tau\).  It solves
\begin{equation}
\max_{w\in\Delta_K,\,z\in\{0,1\}^Q}\sum_{j=1}^Qz_j
\quad\text{subject to}\quad
\langle w,\widehat\Delta_{q_j,a}\rangle
\ge
\varepsilon_N+\tau-B_{\rm M}(1-z_j)
\quad
\forall j,\ a\ne a^\star(q_j).
\label{eq:buffered-milp}
\end{equation}
If \(z_j=1\), every wrong answer must trail the gold answer by at least
\(\varepsilon_N+\tau\); if \(z_j=0\), the big-\(M\) term releases the
constraints.  Thus this is the exact MILP for the \(\tau\)-tightened buffered
objective, and it approaches Eq.~(\ref{eq:buffered-erm}) as
\(\tau\downarrow0\).
The implementation in \texttt{src/tte/milp.py} exposes the threshold as
\texttt{strict\_margin}.  Setting it to \(\varepsilon_N+\tau\) implements the
\(\tau\)-tightened objective.  The proof of Theorem~\ref{thm:buffered-main}
then goes through with the comparator band
\(0<M_q(w^\star)\le2\varepsilon_N+\tau\), provided
\(2\varepsilon_N+\tau\le t_0\).  Existing scripts that use \(10^{-4}\) employ
only a numerical tie-breaking margin and should not be described as evaluations
of the calibrated buffered estimator.
